\documentclass[../numerical_methods.tex]{subfiles}
\begin{document}
\section{Aula 05 - 17 de Março, 2026}
\subsection{Motivações}
\begin{itemize}
	\item Continunado Numpy
\end{itemize}
\subsection{Broadcasting, Comparações e Operações Avançadas}
Para garantir a eficiência computacional na realização de operações entre arrays, o numpy utiliza um mecanismo de extensão de arrays chamado \textbf{broadcasting}, permitindo que as operações aritméticas com as arrays sejam feitas elemento a elemento; a forma como isso é feita é transformando as arrays envolvidas nas operações aritméticas, de modo que tenham as mesmas dimensões, possibilitando justamente a possibilidade de fazer as operações elemento a elemento.
Todas as linhas de código que usarem o numpy devem ser entendidas como tendo o comando \texttt{import numpy as np}.

Quando somamos um escalar a uma array, o numpy primeiro replica o escalar em uma array com a mesma dimensão da array envolvida na operação. Vamos considerar um exemplo na prática para somar uma matriz A com o escalar 3:
\[
	B = s+A = \begin{bmatrix}
		3 & 3 & 3 & 3 & 3 \\
		3 & 3 & 3 & 3 & 3 \\
		3 & 3 & 3 & 3 & 3 \\
		3 & 3 & 3 & 3 & 3 \\
		3 & 3 & 3 & 3 & 3
	\end{bmatrix} + \begin{bmatrix}
		0  & 1  & 2  & 3  & 4  \\
		5  & 6  & 7  & 8  & 9  \\
		10 & 11 & 12 & 13 & 14 \\
		15 & 16 & 17 & 18 & 19 \\
		20 & 21 & 22 & 23 & 24
	\end{bmatrix}
\]
\begin{minted}[bgcolor = DarkBackground, linenos=true, breaklines]{Python}
A = np.arange(25).reshape(5, 5)
s = 3 
B = s + A 
print(B)
\end{minted}

Broadcasting também é aplicado quando operações aritméticas são aplicadas em pares de arrays, mas com algumas restrições:
\begin{itemize}
	\item Os dois arrays devem possuir dimensões compatíveis em algum dos eixos;
	\item O broadcasting é aplicado nos demais exios para que ambos os arrays tenham as mesmas dimensões;
	\item A operação de multiplicação * corresponde a uma multiplicação elemento por elemento, não uma multiplicação matricial.
\end{itemize}
\begin{minted}[bgcolor = DarkBackground, linenos=true]{Python}
A = np.arange(25).reshape(5,5)
v = np.arange(5)
B = v*A
\end{minted}

A ideia do que está acontecendo é:
\begin{align*}
	B = v*A & = \begin{bmatrix}
		            0 & 1 & 2 & 3 & 4 \\
	            \end{bmatrix} * \begin{bmatrix}
		                            0  & 1  & 2  & 3  & 4  \\
		                            5  & 6  & 7  & 8  & 9  \\
		                            10 & 11 & 12 & 13 & 14 \\
		                            15 & 16 & 17 & 18 & 19 \\
		                            20 & 21 & 22 & 23 & 24
	                            \end{bmatrix}               \\
	        & \Rightarrow B = \begin{bmatrix}
		                          0 & 1 & 2 & 3 & 4 \\
		                          0 & 1 & 2 & 3 & 4 \\
		                          0 & 1 & 2 & 3 & 4 \\
		                          0 & 1 & 2 & 3 & 4 \\
		                          0 & 1 & 2 & 3 & 4 \\
	                          \end{bmatrix} * \begin{bmatrix}
		                                          0  & 1  & 2  & 3  & 4  \\
		                                          5  & 6  & 7  & 8  & 9  \\
		                                          10 & 11 & 12 & 13 & 14 \\
		                                          15 & 16 & 17 & 18 & 19 \\
		                                          20 & 21 & 22 & 23 & 24
	                                          \end{bmatrix}
\end{align*}
Quando o array v é em coluna, o broadcasting é feito nelas -- é realmente um tipo de transmissão de informação!
\begin{align*}
	 & B = v*A = \begin{bmatrix}
		             0 \\
		             1 \\
		             2 \\
		             3 \\
		             4
	             \end{bmatrix} * \begin{bmatrix}
		                             0  & 1  & 2  & 3  & 4  \\
		                             5  & 6  & 7  & 8  & 9  \\
		                             10 & 11 & 12 & 13 & 14 \\
		                             15 & 16 & 17 & 18 & 19 \\
		                             20 & 21 & 22 & 23 & 24
	                             \end{bmatrix}       \\
	 & \Rightarrow B = \begin{bmatrix}
		                   0 & 0 & 0 & 0 & 0 \\
		                   1 & 1 & 1 & 1 & 1 \\
		                   2 & 2 & 2 & 2 & 2 \\
		                   3 & 3 & 3 & 3 & 3 \\
		                   4 & 4 & 4 & 4 & 4 \\
	                   \end{bmatrix} * \begin{bmatrix}
		                                   0  & 1  & 2  & 3  & 4  \\
		                                   5  & 6  & 7  & 8  & 9  \\
		                                   10 & 11 & 12 & 13 & 14 \\
		                                   15 & 16 & 17 & 18 & 19 \\
		                                   20 & 21 & 22 & 23 & 24
	                                   \end{bmatrix}
\end{align*}

\begin{minted}[bgcolor = DarkBackground, linenos=true]{Python}
A = np.arange(25).reshape(5,5)
v = np.arange(5).reshape(5,1)
B = v*A
\end{minted}
Outro caso que vale a pena ver é quando os dois operandos são arrays unidimensionais, sendo um deles array linha e outro array coluna; aqui, o broadcasting é aplicado a ambos!
\[
	Z = v+w \Rightarrow Z= \begin{bmatrix}
		0 & 1 & 2 & 3 & 4
	\end{bmatrix} + \begin{bmatrix}
		0 \\
		1 \\
		2
	\end{bmatrix} \Rightarrow Z =\begin{bmatrix}
		0 & 1 & 2 & 3 & 4 \\
		0 & 1 & 2 & 3 & 4 \\
		0 & 1 & 2 & 3 & 4 \\
		0 & 1 & 2 & 3 & 4
	\end{bmatrix}+ \begin{bmatrix}
		0 & 0 & 0 & 0 & 0 \\
		1 & 1 & 1 & 1 & 1 \\
		2 & 2 & 2 & 2 & 2
	\end{bmatrix}
\]
\begin{minted}[bgcolor = DarkBackground, linenos=true]{Python}
v = np.arange(5)
print(v.shape)
w=np.arange(3).reshape(3,1)
print(w.shape)
Z = v+w
print(Z.shape)
\end{minted}
Ele também é usado em atribuições e pode ser forçado para ocorrer em um dos eixos:
\begin{minted}[bgcolor = DarkBackground, linenos=true, breaklines]{Python}
# Usando em atribuições
A = np.zeros((5,3))
A[:3]=[-1,-2,-3]
print(A)

# Forçando o broadcasting em v nas colunas 
v = np.arange(5)
A[:]=v[:, np.newaxis]
print(A)

# O comando acima equivale a fazer 
# A[:]=v.reshape(5,1)
\end{minted}

As comparações usando as operações booleanas já foram vistas na parte das máscaras, mas para verificar se duas arrays são iguais em todos os seus elementos, deve-se utilizar o método \texttt{array\_equal}
\begin{minted}[bgcolor = DarkBackground, linenos=true, breaklines]{Python}
A = np.arange(4).reshape(2,2)
B = np.arange(6,2,-1).reshape(2,2)
C = np.copy(A)

print(np.array_equal(A, B))
print(np.array_equal(A, C))

# Equivalentemente, poderia ser feito 
# D = (A==B)
# print(np.all(D))
\end{minted}

O numpy possui métodos de redução como \texttt{sum}, \texttt{max}, \texttt{min}, \texttt{argmax} e \texttt{argmin}, que podem ser aplicadas à array como um todo ou somente às linhas ou colunas: o que controla a forma de aplicação é o parâmetro \textit{axis}, e quando ele não é especificado, a redução é aplicada à array toda:
\begin{itemize}
	\item \textit{axis não especificado}: retorna um valor único;
	\item \textit{axis=0}: restringe a redução às colunas, percorrendo as linhas, e retorna uma array com número de elementos igual ao número de colunas;
	\item \textit{axis=1}: restringe a redução às linhas, percorrendo as colunas, e retorna uma array com número de elementos igual ao número de linhas.
\end{itemize}
\begin{minted}[bgcolor = DarkBackground, linenos=true, breaklines]{Python}
A = np.zeros((5,5))
A[:]=np.arange(5)

print('Soma de todos os valores: ', np.sum(A))
print('Soma dos valores das colunas: ', np.sum(A, axis=0))
print('Soma dos valores das linhas: ', np.sum(A, axis=1))

B = np.random.randint(low=0, hight=30, size=(5,5))

print('Maior valor dentre todos na matriz: ', np.max(B))
print('Posição do maior valor dentre todos na matriz: ', np.argmax(B))

print('Maior valor em cada colna: ', np.max(B, axis=0))
print('Posição do maior valor em cada coluna: ', np.argmax(B, axis=0))
\end{minted}
As reduções estatísticas, como \texttt{mean}, \texttt{median} e \texttt{std}, funcionam da mesma forma (podem ser aplicadas ao array como um todo, ou especificando linhas ou colunas com o \texttt{axis}), junto aos métodos de reduções lógicas \texttt{all}, verificando se todos os elementos do array satisfazem a condição, e \texttt{any}, que verifica se algum elemento do array satisfaz a condição.

Para ordenar os elementos do array, o método \texttt{sort} é utilizado com a própria variável, ou seja, a ordenação é feita diretamente no array (\textit{inplace}), com um detalhe de que o método sort invocado pelo numpy gera uma cópia ordenada do array --o original não é modificado! Por padrão, a ordenação é feita em cada linha, mas pode ser feita em cada coluna usando o \texttt{axis}.
\begin{minted}[bgcolor = DarkBackground, linenos=true, breaklines]{Python}
A = np.random.randint(low=0,hight=30, size=(5,5))
A.sort()
B=np.sort(A) # Cópia de A, mas ordenada 
A.sort(axis=0) # Ordenando por colunas
\end{minted}

O método \texttt{argsort} retorna um array contendo os índices dos elementos ordenados nas linhas:
\begin{minted}[bgcolor = DarkBackground, linenos=true, breaklines]{Python}
A = np.array([[12.0, 1.0, 26.0],
[23.0, 0.0, 7.0],
[10.0, 19.0, 21.0]])

I = np.argsort(A)
print(I)
for i in range(A.shape[0]):
    print('Índices dos elementos ordenados na linha', i)
    print('Linha', i, ': '+10*' ', A[i])
    print('Linha', i, 'ordenada: ', A[i, I[i]], '\n')
\end{minted}
A multiplicação matricial e vetorial como definida na álgebra matricial é feita utilizando o método \texttt{dot}, mas as dimensões devem ser compatíveis!
\begin{minted}[bgcolor = DarkBackground, linenos=true, breaklines]{Python}
A = np.zeros((2,2))
v = np.array([1,2]).reshape(2,1)
A[:]=v 
b1=np.dot(A,v)
print(b1, b1.shape)

## Atenção ## 
# O produto np.dot(v,A) não é válido, pois as dimensões 
# são incompatíveis.

# Para tornar compatíveis, precisamos transpor o vetor v! 
b2=np.dot(v.T, A)
print(b2, b2.shape)
\end{minted}

Existem muito mais coisas disponíveis apenas com o numpy, incluindo uma biblioteca inteira para resolução de sistema e decomposições em \textit{eigenvectors} (\texttt{numpy.linalg}), outra para aritmética, estimativas e interpolação polinomial (\texttt{numpy.polynomial}) e até transformação de Fourier (\texttt{numpy.fft}).

\end{document}
